\subsection{First variation and resolvent trace}
\label{subsec:first-variation}

For a perturbation $A_g \to A_g + \delta A_g$, the first variation of the determinant satisfies formally
\begin{equation}
  \delta S_\Pi
  =
  \frac12 \Tr\!\left(A_g^{-1}\,\delta A_g\right) .
\end{equation}

This identity shows that the response of the functional is governed by the resolvent kernel.
In three spatial dimensions, this mechanism produces a universal $1/r$ Green profile, leading to a Newtonian response without assuming a fundamental force law.

In the present covariant four-dimensional setting, the metric variation defines a renormalized spectral stress tensor,
\begin{equation}
  \mathcal{T}^{\Pi,{\mathrm ren}}_{\mu\nu}
  \equiv
  -\frac{2}{\sqrt{g}}
  \frac{\delta S_\Pi^{\mathrm ren}}{\delta g^{\mu\nu}} .
\end{equation}

Its precise geometric decomposition depends on the heat-kernel structure of $A_g$.
In particular, local two-derivative, local four-derivative, non-local, and anomaly contributions will appear with distinct origins.
The classification of these terms and their infrared hierarchy is developed in the following sections.
